\section{Ergodic Theory in General State Space}
\label{sec:A.5}

This section shows how to generalize the definitions of irreducibility, aperiodicity and recurrence
to general state spaces.  We will also summarize some ergodicity results in general state space.

\begin{definition}[$\varphi$-Irreducible Markov Chain]
\label{def:A.26}
A Markov chain is called $\varphi$-irreducible if there exists a non-zero measure $\varphi$ on $E$
such that, for all $A \subseteq E$ with $\varphi(A) > 0$ and for all $x \in E$, there exist a
positive integer $n = n(x)$ such that $P^n(x, A) > 0$.
\end{definition}

\begin{definition}[Aperiodic Markov Chain]
\label{def:A.27}
A chain is aperiodic if there do not exist $d \geq 2$ and disjoint subsets $E_1, \ldots, E_d
\subseteq E$, with
\begin{align*}
  P(x, E_{i+1}) &= 1, \quad \forall x \in E_i, \quad i \in \{1, \ldots, d-1\}, \\
  P(x, E_1)     &= 1, \quad \forall x \in E_d.
\end{align*}
\end{definition}

We remark that if $E$ is discrete, this definition agrees with the one previously given.  The
following result ensures that a Markov chain that is $\varphi$-irreducible and aperiodic has a
limit distribution.

\begin{theorem}[Limit Distribution for Almost-All Initial States]
\label{thm:A.28}
The distribution of an aperiodic, $\varphi$-irreducible Markov chain converges to a limit
distribution $\pi$ for almost-all initial states.  More precisely, there is $\pi$ such that, for
$\pi$-almost all $x \in E$,
\[
  \lim_{n \to \infty} d_{\mathrm{TV}}(P^n(x, \cdot), \pi) = 0.
\]
\end{theorem}

The above limit holds for almost all starting values $x \in E$.  To make it hold for \emph{all}
starting values $x \in E$ and thus for all initial distributions $\pi_0$ we need an additional
property: \emph{Harris recurrence}.

\begin{definition}[Harris Recurrence]
\label{def:A.29}
A chain $\mathbb{X}$ is Harris recurrent if there is a probability $\nu$ such that for all
$A \subseteq E$ with $\nu(A) > 0$ and all $x \in E$ we have
\[
  \Prob(X_n \in A \text{ for some } n > 0 \mid X_0 = x) = 1.
\]
\end{definition}

\begin{theorem}[Limit Distribution for All Initial States]
\label{thm:A.30}
The distribution of an aperiodic, Harris recurrent Markov chain converges to a limit distribution
$\pi$, regardless of the initial state.  More precisely, there is a limit distribution $\pi$ such
that, for all $x \in E$,
\[
  \lim_{n \to \infty} d_{\mathrm{TV}}(P^n(x, \cdot), \pi) = 0.
\]
\end{theorem}

\begin{corollary}[Limit Distribution and Ergodicity]
\label{cor:A.31}
Under the conditions of the previous theorem, for all $A \subseteq E$ and all initial conditions
$\pi_0$, it holds that
\[
  \lim_{n \to \infty} \Prob(X_n \in A) = \pi(A).
\]
Therefore, the chain is ergodic.
\end{corollary}

\begin{proof}
The result follows by the dominated convergence theorem, noting that
$\pi_n(A) = \Prob(X_n \in A) = \int_E \pi_0(x) P^n(x, A)\, dx$.
\end{proof}
